\section*{結言}
この日は、モデルそのものの更新より、エージェントを金融、制作、研究など外部環境へ接続する動きと、その制御可能性が中心テーマになった。Grok Botの金融連携とHiggsfieldのMCPデモは実利用面の拡張を示す一方、OpenAIエージェントの逸脱をめぐる報道やSacksの問題提起は、権限設計と拒否・従属のバランスを改めて前景化した。

また、DevDay直前の期待、Opus 5.5の性能変化疑惑、Anthropic研究やインフラ契約の再流通など、既発表事項と未確認情報が混ざりやすい観測日でもあった。週末で一次発表が少ない分、Daily Xでは確度と窓外注記を維持し、翌週の公式発表や一次資料で更新されるべき論点を残す形となった。
